\documentclass[11pt]{article}
\usepackage[utf8]{inputenc}\usepackage[T1]{fontenc}\usepackage{lmodern}
\usepackage{amsmath,amssymb}
\usepackage{graphicx}\usepackage{float}
\usepackage[a4paper,margin=2.2cm]{geometry}
\usepackage{xcolor}\usepackage{listings}\usepackage{mdframed}
\usepackage{enumitem}\usepackage[hidelinks]{hyperref}\usepackage{titlesec}
\definecolor{azul}{HTML}{0D47A1}\definecolor{verde}{HTML}{1B5E20}\definecolor{cinza}{HTML}{555555}
\definecolor{codebg}{HTML}{F4F6F8}\definecolor{cmnt}{HTML}{2E7D32}\definecolor{kw}{HTML}{0D47A1}\definecolor{str}{HTML}{B71C1C}
\definecolor{boxbg}{HTML}{EEF3F8}\definecolor{boxln}{HTML}{0D47A1}
\renewcommand{\contentsname}{Sumário}\renewcommand{\figurename}{Figura}\renewcommand{\tablename}{Tabela}
\titleformat{\section}{\Large\bfseries\color{azul}}{\thesection}{0.6em}{}
\titleformat{\subsection}{\large\bfseries\color{verde}}{\thesubsection}{0.6em}{}
\newmdenv[backgroundcolor=boxbg,linecolor=boxln,linewidth=1pt,roundcorner=4pt,
  innertopmargin=6pt,innerbottommargin=6pt,skipabove=8pt,skipbelow=8pt]{ideia}
\lstset{
 language=Python, basicstyle=\ttfamily\footnotesize, backgroundcolor=\color{codebg},
 keywordstyle=\color{kw}\bfseries, commentstyle=\color{cmnt}\itshape, stringstyle=\color{str},
 showstringspaces=false, breaklines=true, frame=single, framerule=0pt, xleftmargin=5pt,
 aboveskip=6pt, belowskip=6pt, extendedchars=true, inputencoding=utf8,
 literate={á}{{\'a}}1 {é}{{\'e}}1 {í}{{\'i}}1 {ó}{{\'o}}1 {ú}{{\'u}}1
          {â}{{\^a}}1 {ê}{{\^e}}1 {ô}{{\^o}}1 {ã}{{\~a}}1 {õ}{{\~o}}1
          {ç}{{\c c}}1 {à}{{\`a}}1 {Á}{{\'A}}1 {É}{{\'E}}1 {Ã}{{\~A}}1 {Ç}{{\c C}}1
}

\title{\textbf{Guia do Código, Linha a Linha}\\[3pt]
\large Como cada parte funciona --- e como ler cada gráfico}
\author{Material de apoio --- Seminário de Algoritmos Bioinspirados\\
\small Lucas Rivetti \quad Ian Nunes}
\date{Junho de 2026}

\begin{document}
\maketitle
\begin{ideia}
\textbf{Como usar este guia.} Para cada arquivo do projeto, explicamos o que ele
faz, mostramos os trechos de código mais importantes e os traduzimos em
português simples. No fim, há uma seção ensinando \textbf{como interpretar cada
gráfico} gerado. Não é preciso saber programar para acompanhar: foque nos boxes
azuis (a ``intuição'').
\end{ideia}
\tableofcontents
\vspace{4pt}\hrule

% =====================================================================
\section{Visão geral: os arquivos e o fluxo}
O projeto tem 6 módulos Python, cada um com uma responsabilidade:
\begin{center}\small
\begin{tabular}{ll}
\hline
\textbf{Arquivo} & \textbf{Responsabilidade} \\
\hline
\texttt{problema.py} & Define o problema: gera instâncias, avalia energia e desperdício. \\
\texttt{metricas.py} & Dominância, fronteira de Pareto e as métricas (HV, IGD, ONVG, SP). \\
\texttt{vmpacs.py} & O algoritmo de colônia de formigas (VMPACS). \\
\texttt{nsga2.py} & O algoritmo genético NSGA-II (comparação). \\
\texttt{referencia.py} & A fronteira de Pareto ``melhor conhecida'' (referência). \\
\texttt{experimentos.py} & Roda tudo e gera as tabelas e figuras. \\
\hline
\end{tabular}
\end{center}
\textbf{Fluxo:} \texttt{experimentos.py} pede uma instância a \texttt{problema.py},
roda \texttt{vmpacs.py} e \texttt{nsga2.py} sobre ela, mede a qualidade com
\texttt{metricas.py}, calcula a referência com \texttt{referencia.py} e salva
tabelas/figuras.

\begin{ideia}
\textbf{A ``solução'' em uma frase.} Tudo gira em torno de um vetor
\texttt{aloc} de tamanho $n$: \texttt{aloc[i]} é o servidor que recebe a VM
$i$. Por exemplo, \texttt{aloc = [0, 0, 1, 2, 1]} significa: VMs 0 e 1 no
servidor 0; VM 2 no servidor 1; VM 3 no servidor 2; VM 4 no servidor 1.
\end{ideia}

% =====================================================================
\section{\texttt{problema.py} --- o problema}
\subsection{A classe \texttt{Instancia}}
Guarda os dados de um cenário: demandas das VMs, número de servidores, limites e
constantes de energia.
\begin{lstlisting}
class Instancia:
    def __init__(self, Rp, Rm, m, Tp=0.9, Tm=0.9, P_idle=162.0, P_busy=215.0, eps=1e-4):
        self.Rp = np.asarray(Rp)   # demanda de CPU de cada VM
        self.Rm = np.asarray(Rm)   # demanda de memoria de cada VM
        self.n = len(self.Rp)      # numero de VMs
        self.m = m                 # numero de servidores disponiveis
        self.Tp = Tp; self.Tm = Tm # limites de CPU e memoria (90%)
        self.P_idle = P_idle; self.P_busy = P_busy   # potencias (W)
        self.total_cpu = float(np.sum(self.Rp))      # CPU total (constante)
\end{lstlisting}
\begin{ideia}
\texttt{total\_cpu} é a soma das CPUs de todas as VMs. Ela é \textbf{constante}
(todas as VMs rodam) e, como vimos, é a chave para calcular a energia de forma
exata.
\end{ideia}

\subsection{\texttt{gerar\_instancia} --- cria um cenário}
Reproduz o gerador correlacionado do artigo: sorteia CPU e memória de cada VM e,
conforme o parâmetro \texttt{P}, cria correlação entre elas.
\begin{lstlisting}
for i in range(n):
    Rp[i] = rng.random() * (2 * Rp_ref)   # CPU em [0, 2*Rp_ref)
    Rm[i] = rng.random() * Rm_ref         # memoria em [0, Rm_ref)
    r = rng.random()
    if (r < P and Rp[i] >= Rp_ref) or (r >= P and Rp[i] < Rp_ref):
        Rm[i] = Rm[i] + Rm_ref            # cria a correlacao
\end{lstlisting}
\begin{ideia}
\textbf{O que \texttt{P} faz:} \texttt{P}$\approx$0 gera correlação
\emph{negativa} (CPU alta combina com memória baixa); \texttt{P}$\approx$1 gera
correlação \emph{positiva} (alta com alta); \texttt{P}$\approx$0,5, quase nenhuma.
Isso muda o quão difícil é equilibrar os servidores.
\end{ideia}

\subsection{\texttt{avaliar} --- calcula os dois objetivos}
O coração da modelagem. Soma CPU e memória por servidor e calcula energia e
desperdício.
\begin{lstlisting}
np.add.at(usado_cpu, aloc, inst.Rp)   # soma a CPU de cada VM no seu servidor
np.add.at(usado_mem, aloc, inst.Rm)   # idem para memoria
em_uso = usado_cpu > 0                 # servidores que receberam VMs
k = int(np.count_nonzero(em_uso))      # quantos servidores estao ligados

# Objetivo 1 (energia, forma exata): depende so do numero de servidores k
energia = (inst.P_busy - inst.P_idle) * inst.total_cpu + inst.P_idle * k

# Objetivo 2 (desperdicio): soma de |sobra_cpu - sobra_mem| / (usado)
Lp = inst.Tp - usado_cpu[em_uso];  Lm = inst.Tm - usado_mem[em_uso]
W_j = (np.abs(Lp - Lm) + inst.eps) / (usado_cpu[em_uso] + usado_mem[em_uso])
desperdicio = float(np.sum(W_j))
\end{lstlisting}
\begin{ideia}
\texttt{np.add.at(usado\_cpu, aloc, inst.Rp)} é um truque do NumPy: para cada VM
\texttt{i}, soma \texttt{Rp[i]} na posição \texttt{aloc[i]} de \texttt{usado\_cpu}.
No fim, \texttt{usado\_cpu[j]} = CPU total usada no servidor \texttt{j}. A energia
usa a fórmula exata $(P_{busy}-P_{idle})\cdot\text{CPU total}+P_{idle}\cdot k$,
sem ruído numérico. O desperdício é a soma dos $W_j$ dos servidores em uso.
\end{ideia}

\subsection{\texttt{first\_fit} e \texttt{first\_fit\_decreasing}}
Heurística clássica de empacotamento: coloca cada VM no \emph{primeiro} servidor
onde ela cabe; se não couber em nenhum, abre um novo.
\begin{lstlisting}
for i in ordem:                       # percorre as VMs na ordem dada
    for j in range(len(cap_cpu)):     # tenta nos servidores ja abertos
        if cabe(cap_cpu[j], cap_mem[j], inst.Rp[i], inst.Rm[i], inst):
            cap_cpu[j] += inst.Rp[i]; cap_mem[j] += inst.Rm[i]
            aloc[i] = j; colocado = True; break
    if not colocado:                  # nao coube -> abre servidor novo
        aloc[i] = len(cap_cpu); cap_cpu.append(inst.Rp[i]); cap_mem.append(inst.Rm[i])
\end{lstlisting}
O \texttt{first\_fit\_decreasing} (FFD) é o mesmo, mas ordena as VMs da maior
demanda para a menor antes — costuma usar menos servidores. O FFD gera a solução
inicial $S_0$ do VMPACS e serve de \emph{decodificador} do NSGA-II.

% =====================================================================
\section{\texttt{metricas.py} --- comparar fronteiras}
\subsection{\texttt{domina} --- quem é melhor}
\begin{lstlisting}
def domina(a, b):
    # 'a' domina 'b' se nao e pior em nada E e melhor em algo (minimizacao)
    return np.all(a <= b) and np.any(a < b)
\end{lstlisting}
\begin{ideia}
Tradução direta da definição de Pareto: \texttt{a} domina \texttt{b} se
\texttt{a} é $\le$ \texttt{b} em todos os objetivos \emph{e} estritamente menor em
pelo menos um.
\end{ideia}

\subsection{\texttt{frente\_nao\_dominada} --- acha a fronteira}
Marca como ``dominada'' toda solução que alguém domina; sobram as não-dominadas.
\begin{lstlisting}
for i in range(n):
    for j in range(n):
        if i != j and nao_dom[j] and domina(pontos[j], pontos[i]):
            nao_dom[i] = False; break   # alguem domina i -> i sai da fronteira
\end{lstlisting}

\subsection{As quatro métricas}
\begin{itemize}[leftmargin=1.4em,itemsep=2pt]
  \item \texttt{onvg(frente)} = \texttt{len(frente)}: número de soluções
  (maior = melhor).
  \item \texttt{spacing(frente)}: desvio-padrão das distâncias ao vizinho mais
  próximo (menor = mais uniforme).
  \item \texttt{hipervolume(frente, ref)}: área dominada pela fronteira em
  relação a um ponto de referência (maior = melhor). Para 2 objetivos, soma
  faixas:
\begin{lstlisting}
f = f[np.argsort(f[:, 0])]   # ordena por energia
hv = 0.0; prev_y = ref[1]
for x, y in f:
    hv += (ref[0] - x) * (prev_y - y)   # area de cada retangulo
    prev_y = y
\end{lstlisting}
  \item \texttt{igd(frente, referencia)}: distância média da fronteira de
  \emph{referência} até a fronteira obtida (menor = melhor; mede o quão perto a
  solução chegou da ideal).
\end{itemize}

% =====================================================================
\section{\texttt{vmpacs.py} --- a colônia de formigas}
\subsection{\texttt{construir\_solucao} --- uma formiga monta uma alocação}
É o trecho mais importante. A formiga abre um servidor e o enche escolhendo VMs;
quando nada mais cabe, abre o próximo.
\begin{lstlisting}
while disp.any():                 # enquanto houver VMs nao alocadas
    host_cpu = 0.0; host_mem = 0.0 # abre um servidor novo (vazio)
    while True:
        # candidatas: VMs disponiveis que CABEM no servidor atual
        cand = np.where(disp & (host_cpu+Rp <= Tp) & (host_mem+Rm <= Tm))[0]
        if len(cand) == 0: break  # nada mais cabe -> fecha o servidor
        # HEURISTICA: contribuicao parcial aos 2 objetivos (energia e desperdicio)
        eta = eta1 + eta2
        # combina FEROMONIO + HEURISTICA
        score = alpha * tau[cand, host] + (1 - alpha) * eta
        # REGRA: q<=q0 escolhe o melhor (intensifica); senao sorteia (explora)
        if rng.random() <= q0:  escolhida = cand[np.argmax(score)]
        else:                   escolhida = rng.choice(cand, p=score/score.sum())
        # aplica o movimento e a ATUALIZACAO LOCAL do feromonio
        aloc[escolhida] = host; disp[escolhida] = False
        tau[escolhida, host] = (1-rho_l)*tau[escolhida, host] + rho_l*tau0
\end{lstlisting}
\begin{ideia}
Linha a linha: (1) acha as VMs que cabem; (2) calcula a heurística $\eta$ (quanto
cada VM ajuda nos dois objetivos); (3) soma feromônio e heurística no
\texttt{score}; (4) com prob.\ \texttt{q0} pega a de maior \texttt{score}
(intensifica), senão sorteia proporcional (explora); (5) coloca a VM e evapora um
pouco do feromônio daquele movimento (atualização local).
\end{ideia}

\subsection{\texttt{vmpacs} --- o laço principal}
\begin{lstlisting}
S0 = first_fit_decreasing(inst)            # solucao inicial (FFD)
tau0 = 1.0 / (n * (P_S0 + W_S0))           # feromonio inicial
tau = np.full((n, m), tau0)
for t in range(1, M + 1):                  # M iteracoes
    for _ in range(NA):                    # NA formigas constroem solucoes
        aloc = construir_solucao(...)
        atualizar_arquivo(arquivo, objs, aloc, t)   # guarda nao-dominadas
    for item in arquivo:                   # ATUALIZACAO GLOBAL (so as de Pareto)
        reforco = rho_g * lam / (P_S + W_S)
        tau[i, j] = (1-rho_g)*tau[i, j] + reforco
\end{lstlisting}
\begin{ideia}
A cada iteração, todas as formigas constroem soluções; as não-dominadas entram no
\texttt{arquivo} (a fronteira de Pareto); depois, só essas soluções
\emph{reforçam} o feromônio dos seus movimentos. Assim a colônia ``aprende''.
\end{ideia}

% =====================================================================
\section{\texttt{nsga2.py} --- o algoritmo genético}
\subsection{\texttt{ordenacao\_nao\_dominada} --- separa em frentes}
Classifica a população: frente 0 = não-dominados; frente 1 = os dominados só pela
frente 0; e assim por diante.
\begin{lstlisting}
for p in range(n):
    for q in range(n):
        if domina(objs[p], objs[q]): S[p].append(q)   # p domina q
        elif domina(objs[q], objs[p]): cont[p] += 1    # alguem domina p
    if cont[p] == 0: frentes[0].append(p)              # p e nao-dominado
\end{lstlisting}

\subsection{\texttt{distancia\_multidao} --- mantém diversidade}
Mede o quão ``isolada'' cada solução está. As das pontas recebem distância
infinita (são sempre preservadas), as do meio recebem a distância entre os
vizinhos. Isso evita que todas as soluções se amontoem num ponto só.

\subsection{\texttt{crossover\_ox} e \texttt{mutacao\_swap} --- gerar filhos}
Cada indivíduo é uma \textbf{permutação} das VMs (uma ordem). O \emph{Order
Crossover} copia uma fatia de um pai e completa com a ordem do outro; a mutação
troca duas posições. O laço principal (\texttt{nsga2}) repete: gera filhos, junta
com os pais e mantém os \texttt{pop} melhores (por frente e diversidade).

\begin{ideia}
\textbf{Diferença para o VMPACS:} o VMPACS \emph{constrói} soluções com feromônio;
o NSGA-II \emph{evolui} uma população por cruzamento e mutação. Os dois buscam a
mesma fronteira de Pareto, por caminhos diferentes.
\end{ideia}

% =====================================================================
\section{\texttt{referencia.py} --- a fronteira ``melhor conhecida''}
Serve para \emph{visualizar} o trade-off real. Para cada número de servidores
\texttt{k}, distribui as VMs entre os \texttt{k} servidores buscando equilíbrio, e
refina com busca local.
\begin{lstlisting}
for k in range(kmin, kmin + extra + 1):   # varre numeros de servidores
    # varias tentativas: espalha as VMs entre k servidores, balanceando
    a = aloc_k_bins(rng.permutation(n), k, inst)
    # ... guarda a de menor desperdicio e refina com busca_local
    melhor = busca_local(melhor, inst, iters_local, rng)
return pontos_nao_dominados(np.array(solucoes))
\end{lstlisting}
\texttt{busca\_local} tenta mover VMs entre servidores já usados para reduzir o
desperdício total. Variando \texttt{k}, traça-se a curva de compromisso completa.

% =====================================================================
\section{\texttt{experimentos.py} --- rodar tudo}
\begin{lstlisting}
python3 experimentos.py run 0   # roda a correlacao P=0.0 (e salva parcial)
python3 experimentos.py run 1   # P=0.25  ...  ate  run 4 (P=1.0)
python3 experimentos.py agg     # junta tudo: gera tabelas (CSV) e figuras (PNG)
\end{lstlisting}
Para cada correlação, ele roda VMPACS e NSGA-II 10 vezes, calcula as 4 métricas,
e salva. O \texttt{agg} desenha os gráficos de barras e a fronteira de Pareto.

% =====================================================================
\section{Como ler cada gráfico}
Todos os gráficos comparam \textcolor{red}{\textbf{VMPACS (vermelho)}} e
\textbf{NSGA-II (azul-escuro)} nos 5 níveis de correlação \texttt{P} (eixo
horizontal; o número entre parênteses é o coeficiente de correlação aproximado,
de $-0{,}75$ a $+0{,}75$). As barras têm \emph{barras de erro} (o desvio entre as
10 execuções).

\subsection{Gráficos de barras de energia, desperdício e servidores}
\begin{itemize}[leftmargin=1.4em,itemsep=2pt]
  \item \textbf{Energia (W)} e \textbf{Desperdício} e \textbf{nº de Servidores}:
  em todos, \textbf{barra mais baixa = melhor}.
  \item Como ler: para cada \texttt{P}, compare as duas barras. A menor é o
  algoritmo que gastou menos / desperdiçou menos / usou menos servidores naquele
  cenário.
  \item O que esperar: o NSGA-II costuma ter barras um pouco menores (converge
  para soluções de menor energia e desperdício); em $P=0{,}5$ o VMPACS empata ou
  ganha no desperdício.
\end{itemize}

\subsection{Gráficos de Hipervolume e IGD}
\begin{itemize}[leftmargin=1.4em,itemsep=2pt]
  \item \textbf{Hipervolume (HV)}: \textbf{barra mais ALTA = melhor} (cobre mais
  área boa do espaço de objetivos). É a única em que ``maior é melhor''.
  \item \textbf{IGD}: \textbf{barra mais BAIXA = melhor} (chegou mais perto da
  fronteira de referência).
  \item Como ler: HV alto \emph{e} IGD baixo indicam o algoritmo que melhor se
  aproximou do ideal. Nas figuras, o NSGA-II tem HV maior e IGD menor na maioria
  das correlações.
\end{itemize}

\subsection{Gráfico das fronteiras de Pareto}
\begin{ideia}
Este é o gráfico-chave do trabalho. \textbf{Eixos:} horizontal = energia (W);
vertical = desperdício. \textbf{O canto inferior-esquerdo é o melhor} (pouca
energia \emph{e} pouco desperdício).
\end{ideia}
\begin{itemize}[leftmargin=1.4em,itemsep=2pt]
  \item A \textbf{curva cinza} (``fronteira melhor conhecida'') mostra o
  \emph{trade-off real}: cada ponto é um compromisso (menos energia $\to$ mais
  desperdício, e vice-versa). Ela tem vários pontos justamente porque é a curva
  completa de melhores trocas.
  \item Os \textbf{marcadores} (quadrado vermelho = VMPACS; losango azul =
  NSGA-II) mostram \emph{onde cada algoritmo converge}.
  \item Como ler: quanto mais para baixo e para a esquerda, melhor. O NSGA-II
  costuma ``encostar'' na curva (no joelho de menor energia); o VMPACS fica um
  pouco acima dela (minimiza servidores, mas equilibra um pouco menos).
  \item \textbf{Por que a curva é ``rasa''?} Porque, nestas instâncias, os dois
  objetivos quase não conflitam: empacotar bem reduz energia \emph{e} desperdício
  ao mesmo tempo. Por isso o trade-off existe, mas é pequeno — e cada algoritmo
  sozinho converge para perto de um único ponto (o joelho).
\end{itemize}

\paragraph{``Mas não deveria parecer uma exponencial?''} Ótima pergunta --- e a resposta é \emph{sim}: uma fronteira de Pareto costuma ser uma curva decrescente e convexa (formato em ``L''). A nossa \textbf{é} desse formato (Figura~\ref{fig:formag}), com duas particularidades. \textbf{(i) É discreta:} a energia vale $\text{const}+162\cdot k$, com $k$ inteiro igual ao número de servidores; logo a fronteira são \emph{pontos} (de $162$\,W em $162$\,W), não uma curva contínua --- as curvas lisas dos livros vêm de problemas \emph{contínuos}. \textbf{(ii) É curta:} o desperdício \emph{satura} após uns 2--3 servidores a mais que o mínimo (os servidores já ficam equilibrados), então poucos pontos são não-dominados. Vemos, assim, apenas a \emph{parte inicial} daquela curva convexa.

\begin{figure}[h]\centering
\includegraphics[width=0.97\textwidth]{forma_pareto.png}
\caption{A fronteira de Pareto \emph{é} decrescente e convexa (linha vermelha), mas discreta (degraus de $162$\,W por servidor) e curta (o desperdício satura). Pontos cinzas: soluções dominadas; rótulos: número de servidores.}\label{fig:formag}
\end{figure}

\vspace{6pt}
\begin{ideia}
\textbf{Resumo de ``o que é melhor'' em cada gráfico:} energia, desperdício,
servidores, IGD e Spacing $\to$ \textbf{menor é melhor}; ONVG e Hipervolume
$\to$ \textbf{maior é melhor}; nas fronteiras $\to$ \textbf{inferior-esquerda é
melhor}.
\end{ideia}

\end{document}
